\ifdefined\MathNoteBookMode\else
\documentclass[a4paper,11pt]{ctexart}

\input{../mathnote-preamble.tex}

\renewcommand{\notetitle}{常用切线放缩不等式速览}
\renewcommand{\noteauthor}{高中数学导数核心应用}
\renewcommand{\notedate}{\today}

% --- 页面布局 ---
% 使用较小的页边距以容纳更多内容
\geometry{a4paper, left=1.5cm, right=1.5cm, top=1.5cm, bottom=1.5cm}

% --- 图像与 PGFPlots 设置 ---
\usepackage{pgfplots}
\begin{document}
\fi

\pgfplotsset{compat=1.18}
\pgfplotsset{
    miniplot/.style={
        width=5.5cm, height=4.5cm, % 为所有小图设置统一尺寸
        axis lines=middle,
        xtick=\empty, ytick=\empty, % 隐藏刻度
        xlabel={}, ylabel={},
        enlarge x limits=0.1,
        enlarge y limits=0.1,
        samples=50,
        axis line style={-}, % 简洁的坐标轴样式
    }
}

\PageTag[secondary]{切线放缩·单页速览}

% --- 主标题 ---
\begin{center}
  {\sffamily\bfseries\Huge\color{accent}\notetitle\par}
  \vspace{0.3cm}
  {\color{inkgray}\hrulefill\quad\noteauthor\quad\hrulefill\par}
\end{center}


% ==========================================================
%                      指数函数部分
% ==========================================================
\section*{1. 指数函数相关 (Exponential Functions)}

\begin{tcolorbox}[colback=highlightbg, colframe=highlight!75!black, fonttitle=\bfseries, title=\textcolor{white}{$e^x \ge x+1$}]
\begin{minipage}[c]{0.78\textwidth}
    \begin{itemize}
        \item \keyword{成立条件}: $x \in \R$
        \item \keyword{取等条件}: 当且仅当 $x=0$
        \item \keyword{几何意义}: 函数 $e^x$ 的图像恒在其切点 $(0,1)$ 处的切线 $y=x+1$ 的上方（除切点）。
    \end{itemize}
\end{minipage}%
\begin{minipage}[c]{0.25\textwidth}
    \centering
    \begin{tikzpicture}[scale=0.65]
    \begin{axis}[miniplot, xmin=-2.2, xmax=2.2, ymin=-1.2, ymax=4.2]
        \addplot[domain=-2:2, smooth, thick, color=highlight] {exp(x)} node[pos=0.8, right] {$e^x$};
        \addplot[domain=-2:2, thick, color=accent] {x+1} node[pos=0.9, left] {$x+1$};
        \addplot[only marks, mark=*, black] coordinates {(0,1)};
    \end{axis}
    \end{tikzpicture}
\end{minipage}
\end{tcolorbox}

\begin{tcolorbox}[colback=highlightbg, colframe=highlight!75!black, fonttitle=\bfseries, title=\textcolor{white}{$e^x \ge ex$}]
\begin{minipage}[c]{0.78\textwidth}
    \begin{itemize}
        \item \keyword{成立条件}: $x \in \R$
        \item \keyword{取等条件}: 当且仅当 $x=1$
        \item \keyword{几何意义}: 函数 $e^x$ 的图像恒在其切点 $(1,e)$ 处的切线 $y=ex$ 的上方（除切点）。
    \end{itemize}
\end{minipage}%
\begin{minipage}[c]{0.25\textwidth}
    \centering
    \begin{tikzpicture}[scale=0.65]
    \begin{axis}[miniplot, xmin=-1, xmax=3, ymin=-1, ymax=8]
        \addplot[domain=-1:2.5, smooth, thick, color=highlight] {exp(x)} node[pos=0.8, right] {$e^x$};
        \addplot[domain=-0.5:2.5, thick, color=accent] {e*x} node[pos=0.9, above] {$ex$};
        \addplot[only marks, mark=*, black] coordinates {(1,e)};
    \end{axis}
    \end{tikzpicture}
\end{minipage}
\end{tcolorbox}

% ==========================================================
%                      对数函数部分
% ==========================================================
\section*{2. 对数函数相关 (Logarithmic Functions)}

\begin{tcolorbox}[colback=secondarybg, colframe=secondary!60!black, fonttitle=\bfseries, title=\textcolor{white}{$\ln x \le x-1$}]
\begin{minipage}[c]{0.78\textwidth}
    \begin{itemize}
        \item \keyword{成立条件}: $x > 0$
        \item \keyword{取等条件}: 当且仅当 $x=1$
        \item \keyword{几何意义}: 函数 $\ln x$ 的图像恒在其切点 $(1,0)$ 处的切线 $y=x-1$ 的下方（除切点）。
    \end{itemize}
\end{minipage}%
\begin{minipage}[c]{0.25\textwidth}
    \centering
    \begin{tikzpicture}[scale=0.65]
    \begin{axis}[miniplot, xmin=-0.5, xmax=4, ymin=-2, ymax=2.5]
        \addplot[domain=0.1:4, smooth, thick, color=secondary] {ln(x)} node[pos=0.9, above] {$\ln x$};
        \addplot[domain=-0.5:4, thick, color=accent] {x-1} node[pos=0.9, right] {$x-1$};
        \addplot[only marks, mark=*, black] coordinates {(1,0)};
    \end{axis}
    \end{tikzpicture}
\end{minipage}
\end{tcolorbox}

\begin{tcolorbox}[colback=secondarybg, colframe=secondary!60!black, fonttitle=\bfseries, title=\textcolor{white}{$1-\frac{1}{x} \le \ln x$}]
\begin{minipage}[c]{0.78\textwidth}
    \begin{itemize}
        \item \keyword{成立条件}: $x > 0$
        \item \keyword{取等条件}: 当且仅当 $x=1$
        \item \keyword{几何意义}: 函数 $\ln x$ 的图像在函数 $y=1-1/x$ 的图像上方（除 $(1,0)$ 点）。
    \end{itemize}
\end{minipage}%
\begin{minipage}[c]{0.25\textwidth}
    \centering
    \begin{tikzpicture}[scale=0.65]
    \begin{axis}[miniplot, xmin=-0.5, xmax=4, ymin=-3, ymax=2]
        \addplot[domain=0.2:4, smooth, thick, color=secondary] {ln(x)} node[pos=0.9, above] {$\ln x$};
        \addplot[domain=0.2:4, smooth, thick, color=accent] {1-1/x} node[pos=0.8, right] {$1-\frac{1}{x}$};
        \addplot[only marks, mark=*, black] coordinates {(1,0)};
    \end{axis}
    \end{tikzpicture}
\end{minipage}
\end{tcolorbox}

\begin{tcolorbox}[colback=secondarybg, colframe=secondary!60!black, fonttitle=\bfseries, title=\textcolor{white}{$\ln x \le \frac{x}{e}$}]
\begin{minipage}[c]{0.78\textwidth}
    \begin{itemize}
        \item \keyword{成立条件}: $x > 0$
        \item \keyword{取等条件}: 当且仅当 $x=e$
        \item \keyword{几何意义}: 函数 $\ln x$ 的图像恒在过原点且与其相切的直线 $y=x/e$ 的下方（除切点）。
    \end{itemize}
\end{minipage}%
\begin{minipage}[c]{0.25\textwidth}
    \centering
    \begin{tikzpicture}[scale=0.65]
    \begin{axis}[miniplot, xmin=-0.5, xmax=8, ymin=-1, ymax=3]
        \addplot[domain=0.1:8, smooth, thick, color=secondary] {ln(x)} node[pos=0.9, below right] {$\ln x$};
        \addplot[domain=-0.5:8, thick, color=accent] {x/e} node[pos=0.8, above left] {$\frac{x}{e}$};
        \addplot[only marks, mark=*, black] coordinates {(e,1)};
    \end{axis}
    \end{tikzpicture}
\end{minipage}
\end{tcolorbox}

% ==========================================================
%                      三角函数部分
% ==========================================================
\section*{3. 三角函数相关 (Trigonometric Functions)}

\begin{tcolorbox}[colback=accentbg, colframe=accent!75!black, fonttitle=\bfseries, title=\textcolor{white}{$\sin x < x < \tan x$}]
\begin{minipage}[c]{0.78\textwidth}
    \begin{itemize}
        \item \keyword{成立条件}: $x \in (0, \frac{\pi}{2})$
        \item \keyword{取等条件}: 无
        \item \keyword{几何意义}: 在单位圆中，对于同一锐角弧度 $x$，其正弦线长 < 弧长 < 正切线长。
    \end{itemize}
\end{minipage}%
\begin{minipage}[c]{0.25\textwidth}
    \centering
    \begin{tikzpicture}[scale=0.65]
    \begin{axis}[miniplot, xmin=-0.2, xmax=1.6, ymin=-0.2, ymax=4.2, clip=true, samples=100]
        \addplot[domain=0:1.5, smooth, thick, color=secondary] {sin(deg(x))} node[pos=0.7, right] {$\sin x$};
        \addplot[domain=0:1.5, thick, color=accent] {x} node[pos=0.8, left] {$x$};
        \addplot[domain=0:1.35, smooth, thick, color=highlight] {tan(deg(x))} node[pos=0.9, right] {$\tan x$};
    \end{axis}
    \end{tikzpicture}
\end{minipage}
\end{tcolorbox}

\ifdefined\MathNoteBookMode\else
\end{document}
\fi
